Fix a real $c>\tfrac12$. For a prime $p$ and integer $n\ge 0$, define the integer interval
\[
I_{n,p,c}:=\{a\in\mathbb Z: n<a<n+p^c\}.
\]
Question: is it true that for every $c>\tfrac12$ there exists $p_0(c)$ such that for every prime $p\ge p_0(c)$ and every integer $n\ge 0$ there exist $a,b\in I_{n,p,c}$ with
\[
ab\equiv 1\pmod p?
\]
Equivalently: does every length-$\asymp p^c$ interval contain an element whose multiplicative inverse modulo $p$ also lies in the same interval?

Stress points:
- the threshold $c=1/2$ is natural because inversion is a permutation on $(\mathbb Z/p\mathbb Z)^*$ and ``random'' intersection heuristics suggest $p^c$ needs to exceed $\sqrt p$.
